%---------------------------------------------------------------------
% Title page and one orientation page.
%---------------------------------------------------------------------
\thispagestyle{empty}

\begin{tikzpicture}[remember picture, overlay]
  \fill[bmblued] (current page.north west) rectangle
        ([yshift=-12.3cm]current page.north east);
  \fill[bmaccent] ([yshift=-12.3cm]current page.north west) rectangle
        ([yshift=-12.5cm]current page.north east);
\end{tikzpicture}

\vspace*{0.9cm}
{\color{white}
  {\large علامہ اقبال اوپن یونیورسٹی، اسلام آباد}\par
  \vspace{3pt}
  {بی اے \ \textbullet\ \ ایسوسی ایٹ ڈگری \ \textbullet\ \ بی ایس}\par
  \vspace{20pt}
  {\Huge\bfseries مطالعہ قرآنِ حکیم}\par
  \vspace{6pt}
  {\large\bfseries \textenglish{Study of Quran-e-Hakeem}}\par
  \vspace{12pt}
  {\LARGE\bfseries متوقع سوالات اور ان کے جوابات}\par
  \vspace{12pt}
  {\large کورس کوڈ \textenglish{472}}\par
  \vspace{4pt}
  {اٹھارہ سوالات \ \textbullet\ \ چار حقیقی پرچوں سے ماخوذ\par}
}

\vspace{1.0cm}

\begin{tcolorbox}[enhanced, colback=bmtint, colframe=bmblue, boxrule=0.6pt,
                  arc=3pt, left=12pt, right=12pt, top=10pt, bottom=10pt]
\textbf{\color{bmblue}سوالات کہاں سے لیے گئے ہیں؟}\ \
اس کورس کی مقررہ کتاب (\textenglish{318} صفحات) صرف تصویری اسکین ہے --- کوئی
متنی تہہ نہیں، اور جتنی نقل دستیاب ہوئی اس میں باب کی حد بندی بھی واضح نہیں۔
اس لیے تمام سوالات \textbf{چار حقیقی گزشتہ پرچوں} سے لیے گئے ہیں:

\smallskip
\begin{itemize}\setlength{\itemsep}{0pt}\setlength{\topsep}{2pt}
  \item بی اے --- سمسٹر \textbf{بہار \textenglish{2013}} اور \textbf{خزاں
        \textenglish{2013}}
  \item ایسوسی ایٹ ڈگری --- سمسٹر \textbf{خزاں \textenglish{2021}}
  \item بی ایس --- سمسٹر \textbf{بہار \textenglish{2022}}
\end{itemize}

\smallskip
بتیس سوالات یکجا کر کے \textbf{اٹھارہ موضوعات} بنے --- یہی دہراؤ بتاتا ہے کہ
کون سا موضوع سب سے اہم ہے۔

\medskip
\textbf{\color{bmblue}صحتِ مواد۔}\ \
ہر آیت کا حوالہ سورت اور آیت کے نمبر سے دیا گیا ہے اور معیاری مصحف سے ملا کر
لکھا گیا ہے۔ ترجمہ امتحانی مقصد کے لیے سادہ اور مفہومی ہے --- کسی خاص مطبوعہ
ترجمے کا لفظ بہ لفظ دعویٰ نہیں۔

\medskip
\textbf{\color{bmblue}حیثیت۔}\ \
ذاتی مطالعے کے لیے۔ یہ اے آئی او یو کی مطبوعہ کتاب نہیں۔ اسے من و عن نقل کر کے
اسائنمنٹ میں جمع کرانا منع ہے۔
\end{tcolorbox}

\clearpage

%---------------------------------------------------------------------
\section*{پرچے کا خاکہ}
\markboth{پرچے کا خاکہ}{}

پرچے کی ساخت تینوں ادوار میں تقریباً ایک جیسی ہے، مگر ایک تبدیلی نظر انداز
نہیں ہونی چاہیے:

\medskip
\begin{center}
\begin{tabular}{@{}p{4.6cm} p{4.4cm} p{5.0cm}@{}}
\toprule
 & \textbf{بی اے (\textenglish{2013})}
 & \textbf{ایسوسی ایٹ ڈگری / بی ایس} \\
\midrule
کل نمبر            & \textenglish{100} & \textenglish{100} \\
کامیابی کے نمبر     & \textenglish{40}  & \textbf{\textenglish{50}} \\
وقت                & تین گھنٹے          & تین گھنٹے \\
سوالات دیے جاتے ہیں & آٹھ               & آٹھ \\
حل کرنے ہیں         & پانچ              & پانچ \\
فی سوال نمبر        & \textenglish{20}  & \textenglish{20} \\
\bottomrule
\end{tabular}
\end{center}

\medskip
\begin{ghalti}[سب سے مہنگی غلط فہمی]
کامیابی کے نمبر اب \textbf{چالیس نہیں، پچاس} ہیں --- پرانے پرچوں کی بنیاد پر
''چالیس کافی ہیں'' سمجھنا دس نمبر کے فرق سے فیل کرا دیتا ہے۔
\end{ghalti}

\medskip
\noindent
\textbf{\color{bmblue}کون سا موضوع کتنے پرچوں میں آیا}

\smallskip
\begin{center}
\begin{tabular}{@{}p{7.0cm} p{2.2cm} p{4.8cm}@{}}
\toprule
\textbf{موضوع} & \textbf{کتنے پرچوں میں} & \textbf{اس کتابچے میں} \\
\midrule
اہلِ کتاب کی خصوصیات و تحریف       & تین  & سوال \textenglish{3}، \textenglish{14} \\
سود اور قسموں کے احکام             & تین  & سوال \textenglish{8}، \textenglish{17} \\
محکمات، وحدتِ امت، بخل کی مذمت      & دو   & سوال \textenglish{10--11}، \textenglish{14} \\
منافقین کی صفات                     & دو   & سوال \textenglish{2} \\
طلاق، خلع اور حج کے احکام           & دو   & سوال \textenglish{6--7} \\
الفاظ و صیغوں کی پہچان (لازمی نوعیت) & دو   & سوال \textenglish{18} \\
\bottomrule
\end{tabular}
\end{center}

\medskip
\begin{nukta}[قرآن حکیم کے جواب کا سانچہ]
ممتحن \emph{عربی متن، ترجمہ اور حوالہ} تینوں دیکھتا ہے۔ ہر جواب میں:
\textbf{(۱)} آیت \textbf{خوش خط اور اعراب کے ساتھ}، \textbf{(۲)} سلیس
ترجمہ، \textbf{(۳)} حوالہ (سورت و آیت نمبر)، \textbf{(۴)} مختصر تفسیری نکات۔
پوری آیت یاد نہ ہو تو جتنی یاد ہے اتنی \emph{صحیح} لکھیے --- آدھی صحیح آیت
پوری غلط آیت سے بہتر ہے۔
\end{nukta}
